% Options for packages loaded elsewhere
\PassOptionsToPackage{unicode}{hyperref}
\PassOptionsToPackage{hyphens}{url}
\PassOptionsToPackage{dvipsnames,svgnames,x11names}{xcolor}
\PassOptionsToPackage{space}{xeCJK}
\documentclass[
  a4paper,
  10pt]{article}
\usepackage{xcolor}
\usepackage{amsmath,amssymb}
\setcounter{secnumdepth}{5}
\usepackage{iftex}
\ifPDFTeX
  \usepackage[T1]{fontenc}
  \usepackage[utf8]{inputenc}
  \usepackage{textcomp} % provide euro and other symbols
\else % if luatex or xetex
  \usepackage{unicode-math} % this also loads fontspec
  \defaultfontfeatures{Scale=MatchLowercase}
  \defaultfontfeatures[\rmfamily]{Ligatures=TeX,Scale=1}
\fi
\usepackage{lmodern}
\ifPDFTeX\else
  % xetex/luatex font selection
  \setmainfont[]{Times New Roman}
  \setmonofont[]{Menlo}
  \setmathfont[]{STIX Two Math}
  \ifXeTeX
    \usepackage{xeCJK}
    \setCJKmainfont[]{Songti SC}
  \fi
  \ifLuaTeX
    \usepackage[]{luatexja-fontspec}
    \setmainjfont[]{Songti SC}
  \fi
\fi
% Use upquote if available, for straight quotes in verbatim environments
\IfFileExists{upquote.sty}{\usepackage{upquote}}{}
\IfFileExists{microtype.sty}{% use microtype if available
  \usepackage[]{microtype}
  \UseMicrotypeSet[protrusion]{basicmath} % disable protrusion for tt fonts
}{}
\setlength{\emergencystretch}{3em} % prevent overfull lines
\providecommand{\tightlist}{%
  \setlength{\itemsep}{0pt}\setlength{\parskip}{0pt}}
\usepackage[a4paper,top=24mm,bottom=25mm,left=23mm,right=23mm,headheight=15pt]{geometry}
\usepackage{amsmath,amssymb,mathtools,bm}
\usepackage{booktabs,longtable,array}
\usepackage{graphicx}
\usepackage{float}
\usepackage{caption}
\usepackage{enumitem}
\usepackage{fancyhdr}
\usepackage{titlesec}
\usepackage{mdframed}
\usepackage{xcolor}
\usepackage{xurl}
\usepackage{hyperref}
\graphicspath{{../../}}

\definecolor{PaperBlue}{HTML}{000000}
\definecolor{PaperRule}{HTML}{B7C6D5}
\definecolor{PaperShade}{HTML}{F4F7FA}
\definecolor{PaperText}{HTML}{000000}

\hypersetup{
  colorlinks=true,
  linkcolor=PaperBlue,
  urlcolor=PaperBlue,
  citecolor=PaperBlue,
  pdfborder={0 0 0}
}

\color{PaperText}
\linespread{1.08}
\setlength{\parindent}{1.5em}
\setlength{\parskip}{0.25em}
\setlength{\emergencystretch}{3em}
\setlist{itemsep=0.2em,topsep=0.45em}
\setcounter{secnumdepth}{3}
\setcounter{tocdepth}{3}

\titleformat{\section}
  {\Large\bfseries\color{PaperBlue}}
  {\thesection}{0.7em}{}
\titleformat{\subsection}
  {\large\bfseries\color{PaperText}}
  {\thesubsection}{0.65em}{}
\titleformat{\subsubsection}
  {\normalsize\bfseries\color{PaperText}}
  {\thesubsubsection}{0.6em}{}
\titlespacing*{\section}{0pt}{2.2ex plus 0.5ex}{0.9ex}
\titlespacing*{\subsection}{0pt}{1.8ex plus 0.4ex}{0.7ex}
\titlespacing*{\subsubsection}{0pt}{1.5ex plus 0.3ex}{0.5ex}

\captionsetup{
  font=small,
  labelfont=bf,
  textfont=it,
  justification=centering,
  skip=6pt
}

\pagestyle{fancy}
\fancyhf{}
\fancyhead[L]{\small\nouppercase{\leftmark}}
\fancyhead[R]{\small Selene's Blog}
\fancyfoot[C]{\small\thepage}
\renewcommand{\headrulewidth}{0.4pt}
\renewcommand{\footrulewidth}{0pt}
\fancypagestyle{plain}{
  \fancyhf{}
  \fancyhead[R]{\small Selene's Blog}
  \fancyfoot[C]{\small\thepage}
  \renewcommand{\headrulewidth}{0.4pt}
}

\renewenvironment{quote}
  {\begin{mdframed}[
      linewidth=0.7pt,
      linecolor=PaperRule,
      backgroundcolor=PaperShade,
      leftline=true,
      rightline=false,
      topline=false,
      bottomline=false,
      innerleftmargin=10pt,
      innerrightmargin=8pt,
      innertopmargin=7pt,
      innerbottommargin=7pt,
      skipabove=8pt,
      skipbelow=10pt
    ]\itshape}
  {\end{mdframed}}

\newenvironment{theorembox}[1]
  {\par\medskip\noindent\textbf{Theorem (#1).}\quad}
  {\par\medskip}

\newenvironment{lemmabox}[1]
  {\par\medskip\noindent\textbf{Lemma #1.}\quad}
  {\par\medskip}

\newenvironment{proofbox}
  {\par\medskip\noindent\textit{Proof.}\quad}
  {\hfill$\square$\par\medskip}

\AtBeginDocument{\allowdisplaybreaks[2]}
\usepackage{bookmark}
\IfFileExists{xurl.sty}{\usepackage{xurl}}{} % add URL line breaks if available
\urlstyle{same}
% fallback for those not using the hyperref driver hyperxmp:
\makeatletter
\@ifundefined{xmpquote}{\newcommand{\xmpquote}[1]{#1}}{}
\makeatother
\hypersetup{
  pdftitle={HDU Summer League (8)},
  pdfauthor={Selene},
  colorlinks=true,
  linkcolor={black},
  filecolor={Maroon},
  citecolor={Blue},
  urlcolor={black},
  pdfcreator={LaTeX via pandoc}}

\title{HDU Summer League (8)}
\author{Selene}
\date{August 15, 2026}

\begin{document}
\maketitle

\section{HDU Summer League(8)}\label{hdu-summer-league8}

\subsection*{1001}\label{section}
\addcontentsline{toc}{subsection}{1001}

\textbf{Problem}: Given a directed graph with \(n\) vertices and \(m\)
edges, you need to travel from vertex \(1\) to vertex \(n\). There are
\(q\) operations. The \(i\)-th operation deletes all outgoing edges of
vertex \(p_i\). You need to find the maximum number of operations \(k\)
such that after performing the first \(k\) operations, there is still a
path from vertex \(1\) to vertex \(n\).

\textbf{Solution}: First assume that all outgoing edges of every \(p_i\)
have already been deleted. We then add the edges back in reverse order
while maintaining two arrays

\begin{itemize}
\tightlist
\item
  \(on[u]\): whether the outgoing edges of vertex \(u\) currently exist
\item
  \(vis[u]\): whether \(u\) can currently reach \(n\)
\end{itemize}

When restoring vertex \(u=p[i]\) in reverse order, set \(on[u] = 1\). If
there exists an edge \(u \to v\) and \(vis[v] = 1\), then
\(vis[u] = 1\). Then perform a bfs starting from \(u\) on the reverse
graph.

\subsection*{1002}\label{section-1}
\addcontentsline{toc}{subsection}{1002}

\textbf{Problem}: Given two sequences \(a\) and \(b\) of length \(n\),
there are \(m\) days. At the beginning of each day, \(a_i = a_i + b_i\)
is performed, and there is one query every day

\begin{itemize}
\tightlist
\item
  \texttt{1\ l\ r\ x}: for all \(l \le i \le r\), perform
  \(b_i \leftarrow b_i+x\)
\item
  \texttt{2\ l\ r}: calculate \(\sum_{i=l}^{r} a_i \bmod 2^{64}\)
\end{itemize}

\textbf{Solution}: Maintain prefix sums for both arrays. Since the
modulus is \(2^{64}\), we can directly use \texttt{ull}. For the first
operation, use a Fenwick tree to maintain two lazy arrays \texttt{lazy1}
and \texttt{lazy2}. One stores all modifications \(x\), while the other
stores \(day \times x\). On day \(t\), the answer for an interval
\([l,r]\) is

\[
\sum A_i+t\sum B_i+t\sum lazy1_i-\sum lazy2_i
\]

\subsection*{1003}\label{section-2}
\addcontentsline{toc}{subsection}{1003}

\textbf{Problem}: Given a permutation, each query gives \texttt{x\ k}.
We need to select some positions such that \(x\) is not selected and the
distance between any two selected positions is at least \(k\). Find the
number of valid selection schemes.

\textbf{Solution}: We can think about it in reverse and calculate the
total number of schemes minus the number of schemes in which \(x\) is
selected.

Let \(f(n,k)\) denote the number of schemes for \(n\) positions where
the distance between any two selected positions is at least \(k\). For a
query \((x,k)\), the total number of schemes is \(f(n,k)\). For schemes
that must contain \(x\), consider the valid positions on the left and
right of \(x\). The valid length on the left is \(l=\max(0,x-k)\), and
the valid length on the right is \(r=\max(0,n-x-k+1)\). Therefore, the
answer is

\[
f(n,k)-f(l,k)f(r,k)
\]

There are two ways to calculate \(f\):

Recurrence:

\[
f(i,k)=f(i-1,k)+f(i-k,k)
\]

For a fixed \(k\), it can be precomputed in \(O(n)\).

Or enumerate the number \(c\) of selected positions. Transforming the
condition that adjacent selected positions have distance at least \(k\)
into an ordinary position-selection problem gives

\[
f(n,k)=\sum_{c\ge 0}\binom{n-(k-1)(c-1)}{c}
\]

The complexity is \(O(n/k)\).

If there are many queries with the same \(k\), use the first method;
otherwise, calculate it using combinations.

\subsection*{1007}\label{section-3}
\addcontentsline{toc}{subsection}{1007}

\textbf{Problem}: Given an \(n \times m\) grid with obstacles in some
cells, you can move up, down, left, or right from a cell. There are
\(k\) directed portals in the grid, and a cell may contain multiple
portals. You may choose not to use a portal. There are \(q\) queries,
each asking whether it is possible to reach the destination from the
starting point.

\textbf{Solution}: First divide all cells into connected components and
assign an id to the connected component of each cell. Any two cells
inside the same connected component are mutually reachable. Then
consider the portals. A portal \((x_1,y_1) \to (x_2,y_2)\) is equivalent
to adding an edge between the connected components containing these two
cells. We only need to renumber the connected components and build a new
graph, then run bfs from each connected component containing a portal to
determine whether the connected components containing the starting point
and destination are reachable.

\subsection*{1009}\label{section-4}
\addcontentsline{toc}{subsection}{1009}

\textbf{Problem}: Given an array \(P\) of length \(2N\), where \(P\) is
a permutation of \(2N\) and some elements are missing and represented by
\(0\). Define the value of the permutation as

\[
|P_1-P_2|+|P_3-P_4|+\cdots+|P_{2N-1}-P_{2N}|
\]

Find the number of ways to complete the permutation such that its value
is maximized.

\textbf{Solution}: Divide the array into \(N\) pairs of adjacent
positions. If both numbers in a pair are non-zero, its contribution is
already fixed, so we can ignore it. We only consider pairs containing at
least one \(0\). Suppose there are \(k\) such pairs. Sort the
corresponding \(2k\) numbers and divide them into a lower half and an
upper half. To maximize the contribution, each pair must contain one
large number and one small number. Let \(cnt\) be the number of pairs
where both numbers are \(0\). The answer is

\[
l!\times r!\times 2^{cnt}
\]

\subsection*{1012}\label{section-5}
\addcontentsline{toc}{subsection}{1012}

\textbf{Problem}: Given an undirected graph with \(n\) vertices and
\(m\) edges numbered \(1,2,\dots,n\). There may be multiple edges
connecting the same pair of vertices. The input graph is guaranteed to
be bipartite. A cycle is a connected subgraph in which every vertex has
degree \(2\). If a cycle contains exactly \(k\) vertices, it is called a
\(k\)-cycle. The input guarantees that \(k\) is prime. Find the number
of different \(k\)-cycles in the graph.

\textbf{Solution}: Since the graph is bipartite, every cycle must have
even length. Since \(k\) is also prime, the cycle length can only be
\(2\). Therefore, we only need to count the number of ways to choose two
different parallel edges between the same pair of vertices.

\end{document}
